\documentclass{vorlesung}
\usepackage{forest}
\usepackage{german,booktabs}
\usepackage{colortbl}
\usepackage{marvosym,colonequals}
\usepackage[noend,ruled,linesnumbered]{algorithm2e}
\usepackage{listings}             % Include the listings-package
\usepackage{pgfplots}
%\geometry{paper=a4paper} 
\renewcommand*{\proofname}{Beweis}
%\renewcommand*{\corollaryname}{Korollar}
\usetikzlibrary{decorations.pathreplacing,calc,tikzmark,positioning, shapes,fadings,shapes.arrows,shadows,intersections}
\let\checkmark\relax
\usepackage{cancel,dingbat,marvosym,pifont}
\newcommand{\mrk}[1]{\textcolor{ubrot}{\bfseries #1}}
\graphicspath{{img/}}
\uselanguage{german}
\DeclareFontFamily{U}{skulls}{}
\DeclareFontShape{U}{skulls}{m}{n}{ <-> skull }{}
\newcommand{\skull}{\text{\usefont{U}{skulls}{m}{n}\symbol{'101}}}
\newcommand{\ovon}[1]{\ensuremath{\mathcal{O}(#1)}}
\newcommand{\thetavon}[1]{\ensuremath{\Theta(#1)}}
\newcommand{\omegavon}[1]{\ensuremath{\Omega(#1)}}
\SetAlgorithmName{Algorithmus}{algorithm}{Algorithmenverzeichnis}
\deftranslation[to=german]{proof}{Beweis}
\deftranslation[to=german]{Corollary}{Korollar}
\title[\textcolor{black}{AuD}]{Algorithmen \& Datenstrukturen}
%\subtitle{Modul AI-AuD-B\\ \normalsize\hspace*{0.1em} }
\subtitle[\textcolor{black}{Graphalgorithmen}]{Algorithmen auf Graphen III{\\ \normalsize\hspace*{0.1em}Cormen et al., Kapitel 23, 26.1--26.2}}
\date[2023]{12. Einheit, 11.7.2023 {\scriptsize (Revision \gitAuthorIsoDate)}}
\author[Prof.~Wolter u.a.]{Prof.\ Dr.\ Diedrich Wolter -- Smart Environments}
%\date{Sommersemester 2020}
\newcommand{\mycommfont}[1]{\scriptsize\textcolor{ubblau}{#1}}

\tikzfading[name=arrowfading, top color=transparent!0, bottom color=transparent!95]
\tikzset{arrowfill/.style={top color=ubrot!20!yellow, bottom color=ubrot!80, general shadow={fill=black, shadow yshift=-0.8ex, path fading=arrowfading}}}
\tikzset{arrowstyle/.style={draw=ubrot!80,arrowfill, single arrow,minimum height=#1, single arrow,
single arrow head extend=.4cm,}}
\newcommand{\weg}{\stackrel{*}{\to}}
\newcommand{\tikzfancyarrow}[2][2cm]{\tikz[baseline=-0.5ex]\node [arrowstyle=#1] {#2};}

\begin{document}
\lstdefinestyle{customjava}{
  belowcaptionskip=1\baselineskip,
  breaklines=true,
%  frame=L,
  xleftmargin=\parindent,
  language=Java,
  numbers=left,
  showstringspaces=true,
  basicstyle=\scriptsize\ttfamily,
  keywordstyle=\bfseries\color{green!40!black},
  commentstyle=\itshape\color{purple!40!black},
  identifierstyle=\color{blue},
  stringstyle=\color{orange},
}
\lstset{style=customjava}
\begin{frame}
\titlepage
\end{frame}
%
\begin{frame}{Quiz}
\begin{itemize}
%	\item Warum ist topologisches Sortieren ineffizienter als das Sortieren von  Zahlen nach der natürlichen Ordnung (aufsteigend)?
%	\item Wie lässt sich bestimmen, ob ein gerichteter Graph einen Zyklus enthält?
%	\item Welche Komplexität ergibt sich für \textsc{Dijkstra}, wenn er nicht für alle kürzesten Wege von einem Start $s$ zu allen Knoten $v$, sondern nur für ein Paar $u$, $v$ angepasst wird? Wie kann die Anpassung geschehen?
	\item Warum ändert sich nicht die asympototische Zeitkomplexität für \textsc{Dijkstra}, wenn der Algorithmus so modifiziert wird, dass damit alle kürzesten Wege von einem Startknoten (SSSP) bestimmt werden? 
	\item Was ist der Effizienzgewinn von binärer Exponentiation gegenüber iterativer Multiplikation? 
	\item Warum ist relevant, \textsc{Dijkstra} zu betrachten, obwohl der Algorithmus in derselben Komplexitätsklasse \ovon{|V|^3} wie \textsc{Floyd-Warshall} liegt, welcher zugleich APSP löst?
\end{itemize}
\end{frame}
%
\begin{frame}{Inhalt \& Lernziele dieser Einheit}
Das Lernziel dieser Einheit ist, Kenntnisse über das Spektrum graphentheoretischer Methoden zu erweitern, 
sowie ausgewählte Algorithmen auf Graphen \mrk{implementieren und testen} zu können. 
\begin{itemize}
	\item minimale Spannbäume
	\item maximaler Fluss in Netzwerken
\end{itemize}
\end{frame}
%
%
\begin{frame}{Minimale Spannbäume}
\begin{center}
\begin{tikzpicture}[scale=0.55]
	\foreach \x / \y / \l in {-1/0/a, 0/1/b, 2/1/c, 4/1/d, 1/0/e, 3/0/f, 5/0/g, -1/-1.2/h, 4/-1.2/i}{
      \node[fill=ubgelb,circle,draw,inner sep=1pt,minimum size=4mm] (\l) at (\x,\y) {\scriptsize$\l$};
      }
    \foreach \a / \b / \edge / \c in {h/i/{}/8, a/b/{ultra thick}/1, a/e/{ultra thick}/2, b/c/{}/4, c/d/{ultra thick}/1, d/g/{ultra thick}/2, f/g/{ultra thick}/1, e/f/{thick}/5, e/i/{ultra thick}/2, i/g/{ultra thick}/3,a/h/{ultra thick}/2,e/c/{}/4} {
    \draw[\edge] (\a) -- (\b) node[midway,anchor=south,sloped,yshift=-1mm]{\scriptsize $\c$};
    }

 \draw[-] (6,1) -- (6.5,1) node[anchor=west]{\scriptsize Kante };
 \draw[-,ultra thick] (6,0) -- (6.5,0) node[anchor=west]{\scriptsize Kante im Spannbaum};
\end{tikzpicture}
\end{center}
\begin{itemize}
	\item minimale Spannbäume erfassen Zusammenhangsstruktur
	\begin{itemize}
		\item minimale Kantenmenge, die Verbundenheit nicht verändert
		\item Subgraph ist azyklisch $\leadsto$ Spannbaum
	\end{itemize}
	\item Anwendungsbeispiel: Netzwerkanalyse, wichtige Verbindungen identifizieren; Planen notwendiger Verbindungen
	\item Struktur der kürzesten Wege im Graphen
\end{itemize}
\begin{block}{Minimaler-Spannbaum-Problem (engl.\ minimum spanning tree problem)}
Gegeben ein zusammenhängender, ungerichteter Graph $G=(V,E,c)$ mit Kantenbeschriftung $c:E\to \mathbb{N}$. Bestimme $T\subseteq E$, so dass $(V,T)$ ein Baum ist, der $\sum_{\{u, v\}\in T} c(\{u, v\})$ minimiert.
\end{block}
\end{frame}
%
\begin{frame}{Berechnung von Spannbäumen}
\framesubtitle{Grundprinzip}
%\begin{itemize}
%	\item Versuch, einen gierigen Algorithmus zu formulieren
%\end{itemize}
\SetKwInput{KwEingabe}{Eingabe}
\SetKwInput{KwAusgabe}{Ausgabe}
\SetKwProg{myproc}{Procedure}{}{}
\SetKwComment{KwRem}{$\vartriangleright$}{}
\begin{algorithm}[H]
\DontPrintSemicolon
\KwEingabe{zusammenhängender Graph $G=(V,E)$} %, Kantenbeschriftung $c:E\to\mathbb{N}$}
\KwAusgabe{Kanten eines Spannbaums $T\subseteq E$}
$T \gets \emptyset$\;
\While {$\exists k \in E. \{k\}\cup T$ ist zyklenfrei}{%{$\exists v\in V.$  keine Kante $k\in T$ ist inzident mit $v$}{
 Wähle $k \in E$, so dass $\{k\}\cup T$ zyklenfrei bleibt\;
 $T \gets T \cup \{k'\}$\;
}
\caption{\textsc{Extend-Spanning-Tree}}
\end{algorithm}
\begin{theorem}
 Der Algorithmus berechnet einen Spannbaum.
\end{theorem}
\begin{proof}[Beweisskizze]
$T$ bleibt aufgrund der Wahl von $k'$ zyklenfrei. Da $G$ zusammenhängend ist, muss $T$ bei Termination zusammenhängend sein und alle Knoten erreichen ($T$ ist Spannbaum), denn andernfalls gäbe es eine weitere Kante $k'$, die zugefügt werden könnte. 
%Abbruch der While-Schleife, wenn alle Knoten erreicht ($T$ spannt auf), Hinzuname von $k'$ verursacht keinen Zyklus, da $v$ bisher nicht erreichbar. 
%Korrektheit folgt aus Zusammenhang von $G$.
\end{proof}
%\textsc{Extend-Spanning-Tree} ist ein gieriger Algorithmus
\end{frame}
%
\begin{frame}{Berechnung minimaler Spannbäume}
\begin{itemize}
	\item Berechnung minimaler Spannbäume erfordert Auffinden kostenminimalen Kanten, die Verbindung herstellen
	\item Sei $G=(V,E,c)$ ein ungerichteter zusammenhängender Graph mit Kantenbeschriftung $c:E\to\mathbb{N}$
	\begin{itemize}
		\item Eine Partitionierung $V = V_1 \cup V_2, V_1\cap V_2=\emptyset$ heißt \mrk{Schnitt}
		\item Eine $\{u,v\}$ Kante \mrk{kreuzt} einen Schnitt, wenn $u\in V_1, v\in V_2$
		\item Ein Schnitt \mrk{respektiert Kantenmenge $A$}, falls keine Kante aus $A$ den Schnitt kreuzt
		\item Eine  einen Schnitt kreuzende Kante heißt \mrk{leichte Kante}, falls sie die kreuzende Kante mit minimalen Kosten ist
	\end{itemize}
\end{itemize}
\begin{theorem}[Theorem 23.1, Cormen et al.]
Sei $G=(V,E,c)$ ein zusammenhängender ungerichteter Graph mit \label{theo:leicht} Kantenbeschriftung $c:E\to\mathbb{N}$ und $A\subseteq E$ sei Teilmenge eines minimalen Spannbaums von $G$. Wenn $(V_1, V_2)$ ein Schnitt von $G$ ist, der $A$ respektiert, dann ist $A\cup \{a\}$ auch Teilmenge eines minimalen Spannbaumes, falls $a$ eine leichte Kante ist.
\end{theorem}
\end{frame}
%
\begin{frame}{Berechnung minimaler Spannbäume}
\framesubtitle{Eigenschaften}
\begin{itemize}
	\item wichtiges Theorem, denn es erlaubt gierigen Algorithmus
	\begin{itemize}
		\item jede leichte Kante $a$ erweitert minimalen Spannbaum
		\item[$\leadsto$] bestimme iterativ leichte Kanten, vergrößere Kantenmenge $A$
		\item Schleife benötigt $|V|-1$ Iterationen, um alle Knoten zu verbinden
	\end{itemize}
\end{itemize}
\begin{proof}[Beweisskizze]
Beweis durch Widerspruch: Annahme $A\cup\{a\}$ wäre nicht Teil eines minimalen Spannbaum. Dann muss es einen minimalen Spannbaum $T'$ geben, $A\subseteq T'$ mit Kante $a'\in T'$, $a'\neq a$, wobei $a'$ den Schnitt $(V_1, V_2)$ kreuzt.
$T' \cup \{ a \}$ enthält einen Zyklus, denn der Schnitt $(V_1, V_2)$ wird zweimal gekreuzt. Dann ist auch $(T' \setminus \{a'\}) \cup \{a\}$ ein Spannbaum mit kleineren Kosten als $T'$, da $a$ eine leichte Kante ist -- im Widerspruch zur Annahme.
\end{proof}
\end{frame}
%
\begin{frame}{Algorithmus von Kruskal}
\begin{itemize}
	\item Entwickelt von Joseph Kruskal\footnote{Joseph Kruskal (1956). On the shortest spanning subtree and the traveling salesman problem. In: Proceedings of the American Mathematical Society, 7, S.\ 48--50}
	\item \mrk{Idee:} Füge iterativ Teilbäume mit minimalen Kanten zusammen
	\begin{itemize}
		\item Zu Beginn wird zu Graph $G=(V,E)$ als Wald $G_0=(V,\emptyset)$ mit $|V|$ Zusammenhangskomponenten dargestellt
		\begin{itemize}
			\item Jeder Knoten eigene Zusammenhangskomponente
		\end{itemize}
		\item iterativ werden Teilbäume mit leichtester Kante verbunden
	\end{itemize}
\end{itemize}
\begin{center}
\begin{tabular}{lll}
\begin{tikzpicture}[scale=1.0]
	\foreach \x / \y / \l in {0/0/a, 0/1/b, 2/1/c, 1/0/d, 2/0/e}{%
      \node[fill=ubgruen!20,circle,draw,inner sep=1pt,minimum size=4mm] (\l) at (\x,\y) {\scriptsize$\l$};
      }
    \foreach \a / \b / \edge / \c in {a/b/{}/1, b/d/{}/3, a/d/{}/2, b/c/{}/4, c/d/{}/1, e/c/{thick}/2, d/e/{}/3} {
    \draw[\edge] (\a) -- (\b) node[midway,anchor=south,sloped,yshift=-1mm]{\scriptsize $\c$};
}
\end{tikzpicture}
&
\begin{tikzpicture}[scale=1.0]
	\foreach \x / \y / \l in {0/0/a, 0/1/b, 2/1/c, 1/0/d, 2/0/e}{%
      \node[fill=ubgruen!20,circle,draw,inner sep=1pt,minimum size=4mm] (\l) at (\x,\y) {\scriptsize$\l$};
      }
    \foreach \a / \b / \edge / \c in {a/b/{ultra thick}/1, b/d/{}/3, a/d/{}/2, b/c/{}/4, c/d/{}/1, e/c/{}/2, d/e/{}/3} {
    \draw[\edge] (\a) -- (\b) node[midway,anchor=south,sloped,yshift=-1mm]{\scriptsize $\c$};
}
\end{tikzpicture}
&
\begin{tikzpicture}[scale=1.0]
	\foreach \x / \y / \l in {0/0/a, 0/1/b, 2/1/c, 1/0/d, 2/0/e}{%
      \node[fill=ubgruen!20,circle,draw,inner sep=1pt,minimum size=4mm] (\l) at (\x,\y) {\scriptsize$\l$};
      }
    \foreach \a / \b / \edge / \c in {a/b/{ultra thick}/1, b/d/{}/3, a/d/{}/2, b/c/{}/4, c/d/{ultra thick}/1, e/c/{}/2, d/e/{}/3} {
    \draw[\edge] (\a) -- (\b) node[midway,anchor=south,sloped,yshift=-1mm]{\scriptsize $\c$};
}
\end{tikzpicture}
\\
\begin{tikzpicture}[scale=1.0]
	\foreach \x / \y / \l in {0/0/a, 0/1/b, 2/1/c, 1/0/d, 2/0/e}{%
      \node[fill=ubgruen!20,circle,draw,inner sep=1pt,minimum size=4mm] (\l) at (\x,\y) {\scriptsize$\l$};
      }
    \foreach \a / \b / \edge / \c in {a/b/{ultra thick}/1, b/d/{}/3, a/d/{ultra thick}/2, b/c/{}/4, c/d/{ultra thick}/1, e/c/{}/2, d/e/{}/3} {
    \draw[\edge] (\a) -- (\b) node[midway,anchor=south,sloped,yshift=-1mm]{\scriptsize $\c$};
}
\end{tikzpicture}
&
\begin{tikzpicture}[scale=1.0]
	\foreach \x / \y / \l in {0/0/a, 0/1/b, 2/1/c, 1/0/d, 2/0/e}{%
      \node[fill=ubgruen!20,circle,draw,inner sep=1pt,minimum size=4mm] (\l) at (\x,\y) {\scriptsize$\l$};
      }
    \foreach \a / \b / \edge / \c in {a/b/{ultra thick}/1, b/d/{}/3, a/d/{ultra thick}/2, b/c/{}/4, c/d/{ultra thick}/1, e/c/{ultra thick}/2, d/e/{}/3} {
    \draw[\edge] (\a) -- (\b) node[midway,anchor=south,sloped,yshift=-1mm]{\scriptsize $\c$};
}
\end{tikzpicture}
&
\begin{tikzpicture}[scale=1.0]
	\foreach \x / \y / \l in {0/0/a, 0/1/b, 2/1/c, 1/0/d, 2/0/e}{%
      \node[fill=ubgruen!20,circle,draw,inner sep=1pt,minimum size=4mm] (\l) at (\x,\y) {\scriptsize$\l$};
      }
    \foreach \a / \b / \edge / \c in {a/b/{ultra thick}/1,  a/d/{ultra thick}/2,  c/d/{ultra thick}/1, e/c/{ultra thick}/2} {
    \draw[\edge] (\a) -- (\b) node[midway,anchor=south,sloped,yshift=-1mm]{\scriptsize $\c$};
}
\end{tikzpicture}
\end{tabular}
\end{center}
\end{frame}
%
\begin{frame}{Algorithmus von Kruskal}
\SetKwInput{KwEingabe}{Eingabe}
\SetKwInput{KwAusgabe}{Ausgabe}
\SetKwProg{myproc}{Procedure}{}{}
\SetKwComment{KwRem}{$\vartriangleright$}{}
\begin{algorithm}[H]
\DontPrintSemicolon
\KwEingabe{Graph $G=(V,E,c)$, $V=\{v_1, \ldots, v_n\}$ mit Kantenbeschriftung $c:E\to\mathbb{N}$}
\KwAusgabe{Kanten $T$ des minimalen Spannbaums von $G$}
$T\gets \emptyset$\;
\For{$i=1,\ldots n$}{
 $\mathrm{Set}[i] \gets \{v_i\}$\KwRem*[r]{Zusammenhangskomponenten}
}
$E' \gets \text{ sortiere $E$ nach Kosten aufsteigend } c(e), e\in E$\;
\For{$\{v_i, v_j\}\in E'$, in nach Kosten aufsteigender Reihenfolge}{
$i \gets  \textsc{Find-Set}(v_i), j \gets  \textsc{Find-Set}(v_j)$\;
\If {$i \neq j$}{
 $T\gets T \cup \{\{v_i, v_j\}\}$\;
 $\mathrm{Set}[i] \gets \textsc{union}(\mathrm{Set}[i], \mathrm{Set}[j])$\;
 $\mathrm{Set}[j] \gets \emptyset$\KwRem*[r]{nicht mehr benötigt}
}
}
\caption{\textsc{Kruskal}}
\end{algorithm}
\end{frame}
%
\begin{frame}{Algorithmus von Kruskal}
\framesubtitle{Verwaltung der Knotenmengen}
\begin{itemize}
	\item Algorithmus erfordert effizienten Datentyp zum Speichern von Knotenmengen (der Teilbäume)
	\begin{itemize}
	\item Datentyp kann ausnutzen, dass Knotenmengen disjunkt sind
	\end{itemize}
	\item mögliche Implementation auf Basis von Listen
	\begin{itemize}
		\item Disjunktheit erlaubt \textsc{union} durch \textsc{append} zu implementieren
		\item $\mathrm{Set}$ kann also Array von Listen sein
	\end{itemize}
	\item Effizienzsteigerung durch geschickte Implementation\footnote{für Interessierte: siehe Cormen et al., Kapitel 21.1--21.3}
	\begin{itemize}
		\item speichere zusätzlich für jeden Knoten Referenz auf die Menge, in der er enthalten ist $\leadsto$ \textsc{Find-Set} in \ovon{1} 
		\item erweitere Liste um Referenz auf letztes Element $\leadsto$ \textsc{append} in \ovon{1}
		\item speichere Listenlänge, um bei \textsc{union} kürzere Liste anzufügen $\leadsto$ weniger Referenzen müssen aktualisiert werden
	\end{itemize}
\end{itemize}
\end{frame}
%
\begin{frame}{Algorithmus von Kruskal}
\framesubtitle{Analyse}
\begin{theorem}
Algorithmus \textsc{Kruskal} berechnet den minimalen Spannbaum mit \ovon{|E| \log |E|} Zeitaufwand.
\end{theorem}
\begin{itemize}
	\item Zeitaufwand hängt von verwendeter Mengenstruktur ab
	\begin{itemize}
		\item einfache Listen nicht asymptotisch optimal, deshalb hier kein Beweis des Zeitaufwandes
	\end{itemize}
\end{itemize}
\begin{proof}[Beweisskizze, ohne Nachweis der Laufzeit]
Korrektheit folgt direkt aus Theorem von Folie~\pageref{theo:leicht}, da jeweils nur eine leichte Kante ausgewählt wird. Die Laufzeit setzt sich wie folgt zusammen:
\begin{itemize}
	\item Erstellen der $|V|$ Mengenstrukturen $\leadsto$ \ovon{|V|}
	\item Sortieren der Kanten $\leadsto$ \ovon{|E| \log |E|}
	\item \ovon{|E|} Aufrufe von \textsc{Find-Set} und \textsc{union}
\end{itemize}
\end{proof}
\end{frame}
%
\begin{frame}{Flüsse im Netzwerk}
%\framesubtitle{Anwendung von Graphen zur Transportplanung}
\begin{center}
\begin{tikzpicture}[scale=0.65]
	\foreach \x / \y / \l / \col in {0/-1/a/{ubgelb}, 0/1/b/{ubgruen!20}, 2/1/c/{ubgruen!20}, 4/1/d/{ubgelb}, 1/0/e/{ubgruen!20}, 3/-1/f/{ubgruen!20}, 5/0/g/{ubgruen!20}}{
      \node[fill=\col,circle,draw,inner sep=1pt,minimum size=4mm] (\l) at (\x,\y) {\scriptsize$\l$};
      }
    \foreach \a / \b / \edge / \c in {a/b/{}/1, a/e/{}/2, b/c/{}/4, c/d/{}/1, g/d/{}/2, f/g/{}/1, e/f/{}/5, e/c/{}/4} {
    \draw[\edge,->,thick] (\a) -- (\b) node[midway,anchor=south,sloped,yshift=-1mm]{\scriptsize $\c$};
    }
\end{tikzpicture}
\raisebox{10mm}{\tikzfancyarrow{}}
\begin{tikzpicture}[scale=0.65]
	\foreach \x / \y / \l / \col in {0/-1/a/{ubgelb}, 0/1/b/{ubgruen!20}, 2/1/c/{ubgruen!20}, 4/1/d/{ubgelb}, 1/0/e/{ubgruen!20}, 3/-1/f/{ubgruen!20}, 5/0/g/{ubgruen!20}}{
      \node[fill=\col,circle,draw,inner sep=1pt,minimum size=4mm] (\l) at (\x,\y) {\scriptsize$\l$};
      }
    \foreach \a / \b / \edge / \c in {a/b/{}/1, a/e/{}/1, b/c/{}/1, c/d/{}/1, g/d/{}/1, f/g/{}/1, e/f/{}/1} {
    \draw[\edge,->,thick] (\a) -- (\b) node[midway,anchor=south,sloped,yshift=-1mm]{\scriptsize $\c$};
    }
\end{tikzpicture}

\end{center}
\begin{itemize}
	\item Graphen $G=(V,E,c)$ mit Kantenbeschriftung $c:E\to \mathbb{N}$ können auch zur Modellierung von Transportnetzwerken genutzt werden
	\begin{itemize}
	  \item $c(e)$ repräsentiert \mrk{Kapazität} einer Kante, d.h.\ die maximale Menge, die pro Zeiteinheit transportiert werden kann
	  \item Beispiele: Bit/s in Kommunikationsnetzwerken, Kapazität auf Transportrouten eines Logistikunternehmens
	\end{itemize}
	\item \mrk{Aufgabe:} bringe möglichst viel von einer Quelle zu einem Ziel%
	\begin{itemize}
		\item Kapazität von Verbindungen kann nicht überschritten werden
		\item Weglänge wird nicht betrachtet
		\item Beispiel oben: maximaler Fluss von $a$ nach $d$
	\end{itemize}
\end{itemize}
\end{frame}
%
\begin{frame}{Flüsse im Netzwerk}
\framesubtitle{Problemdefinition -- Definitionen}
\begin{itemize}
	\item eine Kantenbeschriftung $c:E\to \mathbb{N}$ in einem gerichteten Graphen heißt \mrk{Kapazitätsfunktion}, $c(e), e\in E$ \mrk{Kapazität}% einer Kante
	\item Es gibt zwei besondere Knoten, $q\in V$, genannt \mrk{Quelle (engl.\ source)}, und $s\in V$, genannt \mrk{Senke (engl.\ sink)}
\end{itemize}
\begin{block}{Flussgraph (engl.\ flow network)}
Ein gerichteter Graph $G=(V,E,c)$ mit Kapazitätsfunktion $c$ mit Quelle $q\in V$ und Ziel $z\in V$ heißt Flussgraph. 
\end{block}
\begin{block}{Fluss im Graphen $G=(V,E,c)$}
Sei $G=(V,E,c)$ ein Flussgraph, dann heißt Funktion $f:V\times V \to \mathbb{R}$ \mrk{Fluss (engl.\ flow)}, wenn sie folgende Bedingungen erfüllt:%
\begin{enumerate}
	\item respektiert Kapazität: $0\leq f(u,v) \leq c(u,v)$ für alle $u,v\in V$
	\item Flussbedingung:$\!\!\sum\limits_{v\in V} f(v,u)\! =\! \sum\limits_{v\in V} f(u,v)$ für alle $u\in V\setminus \{q,s\}$
\end{enumerate}
\end{block}
\end{frame}
%
\begin{frame}{Flüsse im Netzwerk}
\framesubtitle{Problemdefinition}
\begin{block}{Maximaler-Fluss-Problem (engl.\ maximum-flow problem)}
Sei $G=(V,E,c)$ ein Flussgraph mit Quelle $q$. Berechne den maximalen Fluss $f : V\times V \to \mathbb{R}$, d.h.
$$ \sum_{v\in V} f(q,v) \geq \sum_{v\in V} g(q,v)$$
für alle Flüsse $g$.
\end{block}
\begin{itemize}
	\item aufgrund der Flussbedingung gilt Fluss aus Quelle $q$ ist gleich Fluss in Senke $s$
 $$\sum_{v\in V} f(q,v) = \sum_{v\in V} f(v,s)$$
\end{itemize}
\end{frame}
%
\begin{frame}{Flüsse im Netzwerk}
\framesubtitle{Lösungsansatz}
\begin{itemize}
	\item Lösungsansatz von Ford und Fulkerson\footnote{L.R. Ford Jr., D.R Fulkerson (1956). Maximal flow through a network, In: Canadian Journal of Mathematics, 8 S.\ 399--404}
	\begin{itemize}
		\item grundlegend für eine Vielzahl von Flussproblemen und -algorithmen
	\end{itemize}
	\item Ansatz basiert auf mehreren Konzepten
	\begin{itemize}
		\item \mrk{Restnetzwerk}: Flussgraph noch freier Kapazitäten
		\item \mrk{vergrößernder Pfad}: möglicher Fluss entlang Weg $q\weg s$
	\end{itemize}
\end{itemize}
\SetKwInput{KwEingabe}{Eingabe}
\SetKwInput{KwAusgabe}{Ausgabe}
\SetKwProg{myproc}{Procedure}{}{}
\SetKwComment{KwRem}{$\vartriangleright$}{}
\begin{algorithm}[H]
\DontPrintSemicolon
\KwEingabe{Flussgraph $G=(V,E,c)$, Quelle $q$, Senke $s$}
\KwAusgabe{maximaler Fluss $f$}
\For{$u,v \in V$}{
 $f(u,v)\gets 0$\KwRem*[r]{initialisiere Fluss}
 }
 $G_f \gets G$\KwRem*[r]{Restnetzwerk}
 \While{existiert vergrößernder Pfad $p = q\to s$ in $G_f$}{
   erweitere Fluss $f$ um Fluss entlang $p$\;
   passe Kapazitäten in $G_f$ an\;
 }
\caption{\textsc{Ford-Fulkerson}}
\end{algorithm}
\end{frame}
%
\begin{frame}{Restnetzwerk}
\begin{block}{Restnetzwerk (engl.\ residual network)}
Sei $G=(V,E,c)$ ein Flussgraph mit Quelle $q$ und Senke $s$, sowie $f$ ein Fluss in $G$. Dann heißt
$$
c_f(u,v) = \left\{\begin{array}{ll}
c(u,v) - f(u,v) & \text{ falls } (u,v)\in E\\
f(v,u) & \text{ falls } (v,u) \in E\\
0 & \text{ sonst}
\end{array}\right.
$$
die \mrk{Restkapazität} und $G_f=(V,E_,c_f)$ das \mrk{Restnetzwerk} bzgl.\ Fluss $f$ mit $E_f = \{e \in E | c_f(e) > 0\}$.
\end{block}
\begin{itemize}
  \item stellt noch nicht ausgeschöpfte Kapazität dar
  \item Restnetzwerke erlaube Flüsse entgegen gerichteter Kanten
  \begin{itemize}
  	\item Trick der Modellierung erlaubt, Flüsse im iterativen \textsc{Ford-Fulkerson}-Verfahren wieder abzubauen
	\item Flüsse $f(u,v)$ und $f(v,u)$ ergeben zusammen einen effektiven Fluss zwischen $v$ und $u$
  \end{itemize}
\end{itemize}
\end{frame}
%
\begin{frame}{Restnetzwerk}
\framesubtitle{Beispiel}
\begin{center}
\begin{tabular}{ccc}
Flussgraph & Fluss $f$ & Restnetzwerk $G_f$\\[1ex]
\begin{tikzpicture}[scale=1.2]
	\foreach \x / \y / \l / \col in {1/-1/q/{ubgelb}, 0/0/b/{ubgruen!20}, 2/0/c/{ubgruen!20}, 0/1/d/{ubgruen!20}, 2/1/e/{ubgruen!20}, 1/2/s/{ubgelb}}{
      \node[fill=\col,circle,draw,inner sep=1pt,minimum size=4mm] (\l) at (\x,\y) {\scriptsize$\l$};
      }
    \foreach \a / \b / \edge / \c in {q/b/{}/2, q/c/{}/2, b/c/{}/4, b/d/{}/1, c/e/{}/2, e/s/{}/2, d/s/{}/5} {
    \draw[\edge,->,thick] (\a) -- (\b) node[midway,anchor=south,sloped,yshift=-1mm]{\scriptsize $\c$};
    }
\end{tikzpicture}
&
\begin{tikzpicture}[scale=1.2]
	\foreach \x / \y / \l / \col in {1/-1/q/{ubgelb}, 0/0/b/{ubgruen!20}, 2/0/c/{ubgruen!20}, 0/1/d/{ubgruen!20}, 2/1/e/{ubgruen!20}, 1/2/s/{ubgelb}}{
      \node[fill=\col,circle,draw,inner sep=1pt,minimum size=4mm] (\l) at (\x,\y) {\scriptsize$\l$};
      }
    \foreach \a / \b / \edge / \c in {q/b/{ultra thick}/{2}, b/c/{ultra thick}/{2}, c/e/{ultra thick}/{2}, e/s/{ultra thick}/{2}} {
    \draw[\edge,->] (\a) -- (\b) node[midway,anchor=south,sloped,yshift=-1mm]{\scriptsize $\c$};
    }
\end{tikzpicture}
&
\begin{tikzpicture}[scale=1.2]
	\foreach \x / \y / \l / \col in {1/-1/q/{ubgelb}, 0/0/b/{ubgruen!20}, 2/0/c/{ubgruen!20}, 0/1/d/{ubgruen!20}, 2/1/e/{ubgruen!20}, 1/2/s/{ubgelb}}{
      \node[fill=\col,circle,draw,inner sep=1pt,minimum size=4mm] (\l) at (\x,\y) {\scriptsize$\l$};
      }
    \foreach \a / \b / \edge / \c in {q/c/{}/2, b/c/{}/2, b/d/{}/1, d/s/{}/5,s/e/{}/2,e/c/{}/2,b/q/{}/2} {
    \draw[\edge,->,thick] (\a) -- (\b) node[midway,anchor=south,sloped,yshift=-1mm]{\scriptsize $\c$};
    }
    \draw[thick,->,bend right] (c) to (b);
    \node at (1,0.5) {\scriptsize 2};
\end{tikzpicture}
\end{tabular}
\end{center}
\end{frame}
%
\begin{frame}{Vergrößernder Pfad}
\begin{itemize}
  \item Sei $G_r=(V,E,c)$ ein Restnetzwerk mit Weg $p = q\weg s$
  \item Weg $q\weg s$ erlaubt einen Fluss von $q$ nach $s$
  \begin{itemize}
  	\item Kanten in $G_r$ haben positive Kapazität
	\item betrachte $f'$:
	$$f'(u,v) = \left\{\begin{array}{ll} \min\limits_{x\to y\text{ ist Kante in } p} c(x,y) & (u,v)\text{ ist Kante in } p\\
	0 & \text{ sonst }\end{array}\right.
	$$
  \end{itemize}
%  \begin{itemize}
%  \end{itemize}
  	\item Annahme: $G_r$ ist Restnetzwerk von $G$ bzgl.\ eines Flusses $f$
  \item Fluss $f$ kann mit $f'$ vergrößert werden:
  $$
  (f \uparrow f') (u,v) = \left\{\begin{array}{ll}
  f(u,v) + f'(u,v) - f'(v,u) & \text{ falls } (u,v) \in E\\
  0 & \text{ sonst }
  \end{array}\right.
  $$
\end{itemize}
\end{frame}
%
\begin{frame}{Vergrößernde Flüsse}
\framesubtitle{Eigenschaften}
\begin{itemize}
	\item \mrk{Beobachtung:} Flüsse können sich überlagern, sie addieren sich zu einem Gesamtfluss
	\begin{itemize}
		\item \tikz{\node[fill=ubgruen!20,circle,draw,inner sep=1pt,minimum size=4mm] (a) at (0,0) {\scriptsize $a$};\node[fill=ubgruen!20,circle,draw,inner sep=1pt,minimum size=4mm] (b) at (1,0) {\scriptsize $a$}; \draw[ultra thick, ->] (a) -- (b) node[midway,above]{\scriptsize $x/y$};} bezeichnet Fluss $f(a,b)=x$ bei Kapazität $y$, Kapazität im Restnetzwerk $c_f(a,b) = y-x$, $c_f(b,a)=x$
	\end{itemize}
\end{itemize}

\begin{center}
\begin{tabular}{cccc}
Flussgraph & Fluss $f$ & Fluss $f'$ & $f \uparrow f'$\\
\begin{tikzpicture}[scale=0.8]
	\foreach \x / \y / \l / \col in {1/-1/a/{ubgelb}, 0/0/b/{ubgruen!20}, 2/0/c/{ubgruen!20}, 0/1/d/{ubgelb}, 2/1/e/{ubgruen!20}, 1/2/f/{ubgruen!20}}{
      \node[fill=\col,circle,draw,inner sep=1pt,minimum size=4mm] (\l) at (\x,\y) {\scriptsize$\l$};
      }
    \foreach \a / \b / \edge / \c in {a/b/{}/2, a/c/{}/2, b/c/{}/4, b/d/{}/1, c/e/{}/2, e/f/{}/2, d/f/{}/5} {
    \draw[\edge,->] (\a) -- (\b) node[midway,anchor=south,sloped,yshift=-1mm]{\scriptsize $\c$};
    }
\end{tikzpicture}
&
\begin{tikzpicture}[scale=0.8]
	\foreach \x / \y / \l / \col in {1/-1/a/{ubgelb}, 0/0/b/{ubgruen!20}, 2/0/c/{ubgruen!20}, 0/1/d/{ubgelb}, 2/1/e/{ubgruen!20}, 1/2/f/{ubgruen!20}}{
      \node[fill=\col,circle,draw,inner sep=1pt,minimum size=4mm] (\l) at (\x,\y) {\scriptsize$\l$};
      }
    \foreach \a / \b / \edge / \c in {a/b/{ultra thick}/{2/2}, a/c/{}/2, b/c/{ultra thick}/{2/4}, b/d/{}/1, c/e/{ultra thick}/{2/2}, e/f/{ultra thick}/{2/2}, d/f/{}/5} {
    \draw[\edge,->] (\a) -- (\b) node[midway,anchor=south,sloped,yshift=-1mm]{\scriptsize $\c$};
    }
\end{tikzpicture}
&
\begin{tikzpicture}[scale=0.8]
	\foreach \x / \y / \l / \col in {1/-1/a/{ubgelb}, 0/0/b/{ubgruen!20}, 2/0/c/{ubgruen!20}, 0/1/d/{ubgelb}, 2/1/e/{ubgruen!20}, 1/2/f/{ubgruen!20}}{
      \node[fill=\col,circle,draw,inner sep=1pt,minimum size=4mm] (\l) at (\x,\y) {\scriptsize$\l$};
      }
    \foreach \a / \b / \edge / \c in {a/b/{}/{2/2}, a/c/{ultra thick}/{1/2}, b/c/{}/{2/4}, b/d/{ultra thick}/{1/1}, c/e/{}/{2/2}, e/f/{}/{2/2}, d/f/{ultra thick}/{1/5}} {
    \draw[\edge,->] (\a) -- (\b) node[midway,anchor=south,sloped,yshift=-1mm]{\scriptsize $\c$};
    }
    \draw[ultra thick,->,in=45,out=135] (c) to (b);
    \node at (1.0,0.75) {\scriptsize\bfseries{1/2}};
\end{tikzpicture}
&
\begin{tikzpicture}[scale=0.8]
	\foreach \x / \y / \l / \col in {1/-1/a/{ubgelb}, 0/0/b/{ubgruen!20}, 2/0/c/{ubgruen!20}, 0/1/d/{ubgelb}, 2/1/e/{ubgruen!20}, 1/2/f/{ubgruen!20}}{
      \node[fill=\col,circle,draw,inner sep=1pt,minimum size=4mm] (\l) at (\x,\y) {\scriptsize$\l$};
      }
    \foreach \a / \b / \edge / \c in {a/b/{}/{2/2}, a/c/{}/{1/2}, b/c/{}/{1/4}, b/d/{}/{1/1}, c/e/{}/{2/2}, e/f/{}/{2/2}, d/f/{}/{1/5}} {
    \draw[\edge,->] (\a) -- (\b) node[midway,anchor=south,sloped,yshift=-1mm]{\scriptsize $\c$};
    }
%    \draw[ultra thick,->,in=45,out=135] (c) to (b);
%    \node at (1.0,0.75) {\scriptsize\bfseries{1/2}};
\end{tikzpicture}
\end{tabular}
\end{center}
\end{frame}
%
\begin{frame}{Ford-Fulkerson}
\framesubtitle{Analyse}
\begin{itemize}\vspace*{-3mm}
	\item \mrk{Fragen:} 
	\begin{itemize}
	  \item Warum terminiert \textsc{Ford-Fulkerson}?
	  \item Warum ergibt sich der maximale Fluss?
	  \item Welcher vergrößernde Pfad soll ausgewählt werden?
	  \item Welche Laufzeit hat das Verfahren?
	 \end{itemize}
\end{itemize}\vspace*{-3mm}
\begin{theorem}
\textsc{Ford-Fulkerson} angewendet auf einen Flussgraphen $G=(V,E,c)$ mit Quelle $q$ und Senke $s$ terminiert mit dem maximalen Fluss in \ovon{|f^*| \cdot |E|} Zeit, wobei $f^*$ der maximale Fluss ist und $|f^*| = \sum_{v\in V} f^*(q,v) - f^*(v,q)$.
\end{theorem}
\begin{itemize}\vspace*{-2mm}
	\item Theorem besagt, dass jeder vergrößender Fluss gewählt werden darf, um maximalen Fluss zu erhalten
	\begin{itemize}
		\item keine diesbezügliche Strategie in \textsc{Ford-Fulkerson}
		\item Wahl des kürzesten Weges ($\leadsto$ Breitensuche) verringert Laufzeit auf \ovon{|V|\cdot |E|^2}, bekannt als Edmonds-Karp-Algorithmus\footnote{\scriptsize J.\ Edmonds und R.M.\ Karp (1972). Theoretical Improvements in Algorithmic Efficiency for Network Flow Problems. In: Journal of the ACM, 19(2), S.\ 248--264}
	\end{itemize}
\end{itemize}
\end{frame}
%
\begin{frame}{Ford-Fulkerson}
\framesubtitle{Analyse}
\begin{proof}[Beweisskizze, Details siehe Cormen et al. Kapitel 26.2]
Kern des Beweises ist das sog.\ Maximalfluss-Minimalschnitt-Theorem, welches Äquivalenz folgender Aussagen beweist:
\begin{enumerate}
	\item $f$ ist Maximalfluss in $G$
	\item Im Restnetzwerk $G_f$ gibt es keine vergrößernden Wege
	\item $|f| = c(Q,S)$ für einen Schnitt $(Q,S)$ von $G$
\end{enumerate}
Man betrachtet dazu Schnitte, die Knotenmenge $Q$ mit $q\in Q$ und Knotenmenge $S$, $s\in S$ verbinden.
Die Laufzeit ergibt sich aus der Überlegung, wie viele Iterationen der Schleife -- Pfad suchen, Datenstrukturen aktualisieren -- notwendig sind; ein einzelner Schleifenduchlauf ist in \ovon{|E|} möglich.
\end{proof}
\vspace*{-3mm}
\begin{center}
\begin{tikzpicture}[scale=1.0]
	\foreach \x / \y / \l / \col in {0/0.5/q/{ubgelb}, 2/0/b/{ubgruen!20}, 2/1/c/{ubgruen!20}, 4/0/d/{ubgruen!20}, 4/1/e/{ubgruen!20}, 6/0.5/s/{ubgelb}}{
      \node[fill=\col,circle,draw,inner sep=1pt,minimum size=4mm] (\l) at (\x,\y) {\scriptsize$\l$};
      }
    \foreach \a / \b / \edge / \c in {q/c/{}/{11/16}, q/b/{}/{8/13}, c/e/{}/{12/12}, e/s/{}/{15/20}, d/s/{}/{4/4}, b/d/{}/{11/14}, e/b/{}/{4/9}, b/c/{}/{1/4}, d/e/{}/{7/7}} {
    \draw[\edge,->] (\a) -- (\b) node[midway,anchor=south,sloped,yshift=-1mm]{\scriptsize $\c$};
    }
\draw[ubrot] (-0.5,-0.25) rectangle (2.5,1.25);
\draw[ubrot,fill,opacity=0.1] (-0.5,-0.25) rectangle (2.5,1.25);
\node[ubrot,anchor=north west] at (-0.5,1.25) {\scriptsize $Q$};
\draw[ubblau] (3.5,-0.25) rectangle (6.5,1.25);
\draw[ubblau,fill,opacity=0.1] (3.5,-0.25) rectangle (6.5,1.25);\node[ubblau,anchor=north east] at (6.5,1.25) {\scriptsize $S$};
\node[anchor=north west,text width=3cm] at (6.5,1.5){\scriptsize \begin{eqnarray*}c(Q,S)&=&26\\ |f| &=& 19 \end{eqnarray*}\par};
\end{tikzpicture}
\end{center}
\end{frame}
%
\begin{frame}{Zusammenfassung}
\framesubtitle{Algorithmen auf Graphen}
\begin{itemize}
	\item Graphen erlauben, Eigenschaften eines Netzwerkes zu modellieren
	\begin{itemize}
		\item Wege, Transportverbindungen
		\item abstrakte Beziehungen von Entitäten, Relationen 
		\item Varianten von Graphen (z.B.\ durch Beschriftungen) für spezielle Modellierungsaufgaben
	\end{itemize}
%	\item Graphen verallgemeinern Bäume
	\item Varianten der \mrk{Suche von Wegen} in Graphen
%	\begin{itemize}
%		\item wenigste Kanten $\leadsto$ Breitensuche
%		\item kürzester Weg für ein Paar $u,v$ (SPSP), kürzester Weg für einen Startpunkt zu allen Knoten (SSSP) $\leadsto$ \textsc{Dijkstra}
%		\item alle kürzesten Wege (ASSP) $\leadsto$ \textsc{Floyd-Warshall}
%		\item  Reduktionsbeziehung SPSP $\leq_T$ SSSP $\leq_T$ APSP
%%	\begin{itemize}
%%	\end{itemize}
%	\end{itemize}
		\item Grundlage für graphenbasierte Optimierungsprobleme
	\item \mrk{minimale Spannbäume} $\leadsto$ \textsc{Kruskal}
%	\begin{itemize}
%		\item algorithmisches Prinzip der \mrk{dynamischen Programmierung}
%	\end{itemize}
	\item \mrk{maximaler Fluss} in Netzwerken $\leadsto$ \textsc{Ford-Fulkerson}
\end{itemize}
\end{frame}
\end{document}

\end

\pi_i,j Vorgänger auf Weg von Knoten i zu Knoten j




